% =========================================================
%  ppt/06-ai-npc-chase.tex
%  Materi 6 — Pengantar AI NPC dan Perilaku Chase
%  Bab VIII buku ajar
% =========================================================
\documentclass[aspectratio=169,10pt]{beamer}
% =========================================================
%  ppt/config/packages.tex
%  Paket bersama untuk seluruh presentasi kuliah
%  Pengantar Pemrograman Game
% =========================================================

\usepackage[T1]{fontenc}
\usepackage[utf8]{inputenc}
\usepackage[indonesian]{babel}
\usepackage{lmodern}
\usepackage{textcomp}
\usepackage{amsmath}
\usepackage{amssymb}
\usepackage{graphicx}
\usepackage{booktabs}
\usepackage{tabularx}
\usepackage{array}
\usepackage{multicol}
\usepackage{listings}
\usepackage{upquote}
\usepackage{adjustbox}
\usepackage{tcolorbox}
\usepackage{tikz}
\usetikzlibrary{
  arrows.meta,
  positioning,
  shapes.geometric,
  shapes.misc,
  calc,
  fit,
  backgrounds,
  matrix,
  chains,
  decorations.pathreplacing,
  patterns
}

\tcbuselibrary{skins}

% Jangan memuat geometry: ukuran halaman diatur oleh beamer.

% =========================================================
%  ppt/config/theme.tex
%  Satu tema Beamer untuk keempat belas presentasi
%  Pengantar Pemrograman Game
% =========================================================

% --------------------------------------------------------- warna
\definecolor{ubNavy}{HTML}{1F3864}
\definecolor{ubAccent}{HTML}{2E75B6}
\definecolor{ubHead}{HTML}{D9E2F3}
\definecolor{ubSoft}{HTML}{F2F6FC}
\definecolor{ubLine}{HTML}{BFBFBF}
\definecolor{ubText}{HTML}{1F2937}
\definecolor{ubMuted}{HTML}{5A6678}
\definecolor{ubGreen}{HTML}{1E7B45}
\definecolor{ubGreenBg}{HTML}{F1F8F2}
\definecolor{ubRed}{HTML}{B3261E}
\definecolor{ubRedBg}{HTML}{FDF2F1}
\definecolor{ubWarn}{HTML}{A35A05}
\definecolor{ubWarnBg}{HTML}{FFF8EC}
\definecolor{ubGray}{HTML}{F6F6F6}
\definecolor{ubCodeBg}{HTML}{10233F}
\definecolor{ubCodeText}{HTML}{E8EEF7}
\definecolor{ubCodeKey}{HTML}{8FBCF0}
\definecolor{ubCodeCom}{HTML}{8FB37A}
\definecolor{ubCodeStr}{HTML}{D9A46B}
\definecolor{ubCodeNum}{HTML}{7E99BE}

% Nama warna yang dipakai gambar TikZ buku ajar.
\definecolor{ubPrimary}{HTML}{1F3864}
\definecolor{ubAlt}{HTML}{F2F6FC}

% --------------------------------------------------------- dasar
\usetheme{default}
\usefonttheme{professionalfonts}
\setbeamertemplate{navigation symbols}{}
\setbeamerfont{frametitle}{size=\Large,series=\bfseries}
\setbeamerfont{framesubtitle}{size=\small}
\setbeamerfont{footline}{size=\tiny}
\setbeamerfont{block title}{size=\small,series=\bfseries}
\setbeamerfont{block body}{size=\small}
\setbeamerfont{itemize/enumerate body}{size=\small}
\setbeamerfont{itemize/enumerate subbody}{size=\footnotesize}

\setbeamersize{text margin left=0.60cm,text margin right=0.60cm}
\uchyph=0

\setbeamercolor{normal text}{fg=ubText,bg=white}
\setbeamercolor{structure}{fg=ubAccent}
\setbeamercolor{frametitle}{fg=ubNavy,bg=white}
\setbeamercolor{title}{fg=ubNavy}
\setbeamercolor{footline}{fg=ubMuted,bg=white}
\setbeamercolor{itemize item}{fg=ubAccent}
\setbeamercolor{itemize subitem}{fg=ubAccent}
\setbeamercolor{block title}{fg=white,bg=ubNavy}
\setbeamercolor{block body}{fg=ubText,bg=ubSoft}
\setbeamercolor{block title alerted}{fg=white,bg=ubWarn}
\setbeamercolor{block body alerted}{fg=ubText,bg=ubWarnBg}
\setbeamercolor{block title example}{fg=white,bg=ubGreen}
\setbeamercolor{block body example}{fg=ubText,bg=ubGreenBg}
\setbeamercolor{caption name}{fg=ubMuted}

\setbeamertemplate{blocks}[rounded][shadow=false]
\setbeamertemplate{itemize item}{\raisebox{0.10em}{\color{ubAccent}\scriptsize$\blacksquare$}}
\setbeamertemplate{itemize subitem}{\raisebox{0.06em}{\color{ubAccent}\tiny$\blacktriangleright$}}
\setbeamertemplate{enumerate item}{%
  \makebox[1.3em][l]{\color{ubAccent}\insertenumlabel.}}

\addtobeamertemplate{itemize/enumerate body begin}{%
  \setlength{\itemsep}{0.42em}\setlength{\parskip}{0pt}}{}

% --------------------------------------------------------- judul frame
\setbeamertemplate{frametitle}[default][left]

% --------------------------------------------------------- kaki
\newcommand{\materino}{00}
\newcommand{\materifoot}{Pengantar Pemrograman Game \textperiodcentered\ Materi \materino}

\setbeamertemplate{footline}{%
  \begin{beamercolorbox}[wd=\paperwidth,ht=0.015cm,dp=0cm]{footline}%
    \color{ubHead}\rule{\paperwidth}{0.4pt}%
  \end{beamercolorbox}%
  \begin{beamercolorbox}[wd=\paperwidth,ht=0.40cm,dp=0.14cm,%
                         leftskip=0.60cm,rightskip=0.60cm]{footline}%
    \usebeamerfont{footline}\color{ubMuted}\materifoot\hfill\insertframenumber%
  \end{beamercolorbox}%
}

% --------------------------------------------------------- makro buku ajar
\providecommand{\code}[1]{\texttt{\detokenize{#1}}}
\providecommand{\istilah}[1]{\emph{#1}}
\providecommand{\CPMK}[1]{CPMK\,#1}
\providecommand{\CPL}[1]{CPL\,#1}
\providecommand{\SUBCPMK}[1]{Sub-CPMK\,#1}

% --------------------------------------------------------- gaya kode
\lstdefinelanguage{CSharp}{
  morekeywords={abstract,as,base,bool,break,byte,case,catch,char,class,const,
    continue,decimal,default,delegate,do,double,else,enum,event,explicit,
    extern,false,finally,fixed,float,for,foreach,goto,if,implicit,in,int,
    interface,internal,is,lock,long,namespace,new,null,object,operator,out,
    override,params,private,protected,public,readonly,ref,return,sbyte,
    sealed,short,sizeof,stackalloc,static,string,struct,switch,this,throw,
    true,try,typeof,uint,ulong,unchecked,unsafe,ushort,using,virtual,void,
    volatile,while,var,get,set},
  sensitive=true,
  morecomment=[l]{//},
  morecomment=[s]{/*}{*/},
  morestring=[b]",
}

\lstdefinestyle{kode}{
  basicstyle=\ttfamily\footnotesize\color{ubCodeText},
  keywordstyle=\color{ubCodeKey}\bfseries,
  commentstyle=\color{ubCodeCom}\itshape,
  stringstyle=\color{ubCodeStr},
  numberstyle=\tiny\color{ubCodeNum},
  numbers=left,
  numbersep=7pt,
  backgroundcolor=\color{ubCodeBg},
  frame=none,
  showstringspaces=false,
  breaklines=false,
  tabsize=4,
  columns=fullflexible,
  keepspaces=true,
  xleftmargin=12pt,
  xrightmargin=2pt,
  aboveskip=0pt,
  belowskip=0pt,
  upquote=true,
}
\lstset{style=kode}

\newcommand{\kodejelas}[1]{%
  \vspace{0.16cm}%
  {\footnotesize\color{ubMuted}#1\par}%
}
\newcommand{\sumber}[1]{%
  \vspace{0.08cm}%
  {\scriptsize\color{ubMuted}\itshape #1\par}%
}
\newcommand{\awalan}[1]{%
  {\small\color{ubMuted}#1\par}\vspace{0.18cm}%
}

% --------------------------------------------------------- kotak
\newtcolorbox{kotak}[1][]{enhanced, colback=ubSoft, colframe=ubAccent,
  boxrule=0pt, leftrule=2.2pt, arc=1pt, left=6pt, right=6pt, top=4pt,
  bottom=4pt, fontupper=\footnotesize, #1}
\newtcolorbox{peringatan}[1][]{enhanced, colback=ubWarnBg, colframe=ubWarn,
  boxrule=0pt, leftrule=2.2pt, arc=1pt, left=6pt, right=6pt, top=4pt,
  bottom=4pt, fontupper=\footnotesize, #1}
\newtcolorbox{kiat}[1][]{enhanced, colback=ubGreenBg, colframe=ubGreen,
  boxrule=0pt, leftrule=2.2pt, arc=1pt, left=6pt, right=6pt, top=4pt,
  bottom=4pt, fontupper=\footnotesize, #1}

% --------------------------------------------------------- sampul
% \deckcover{nomor}{minggu}{judul}{sub-cpmk}{teknologi}{sumber}{blok}
\newcommand{\deckcover}[7]{%
\begin{frame}[plain]
\begin{tikzpicture}[remember picture, overlay]
  \fill[ubNavy] (current page.south west) rectangle ([xshift=6.00cm]current page.north west);
  \fill[ubAccent] ([xshift=6.00cm]current page.south west) rectangle ([xshift=6.08cm]current page.north west);

  \node[anchor=north west, text width=4.72cm, align=left, inner sep=0pt]
    at ($(current page.north west)+(0.66cm,-0.60cm)$) {%
      {\color{white}\bfseries\fontsize{10.2}{12.2}\selectfont
       Universitas Bale Bandung}\\[4pt]
      {\color{ubHead}\fontsize{7.6}{9.4}\selectfont
       Fakultas Teknologi Informasi}\\[1pt]
      {\color{ubHead}\fontsize{7.6}{9.4}\selectfont
       Program Studi Teknik Informatika}\\[7pt]
      {\color{ubAccent}\rule{4.0cm}{0.5pt}}%
    };

  \node[anchor=north west, text width=4.72cm, align=left, inner sep=0pt]
    at ($(current page.north west)+(0.66cm,-2.72cm)$) {%
      {\color{ubAccent}\fontsize{6.4}{7.8}\selectfont MATA KULIAH}\\[3pt]
      {\color{white}\bfseries\fontsize{9.8}{12.0}\selectfont
       Pengantar Pemrograman Game}\\[3pt]
      {\color{ubHead}\fontsize{7.6}{9.4}\selectfont
       MWCA009 \textperiodcentered\ 3 SKS}\\[8pt]
      {\color{ubHead}\fontsize{7.6}{9.4}\selectfont
       Semester Ganjil 2026/2027}%
    };

  \node[anchor=south west, text width=4.72cm, align=left, inner sep=0pt]
    at ($(current page.south west)+(0.66cm,0.52cm)$) {%
      {\color{ubAccent}\fontsize{6.4}{7.8}\selectfont DOSEN PENGAMPU}\\[3pt]
      {\color{white}\fontsize{8.4}{10.4}\selectfont
       Cecep Suwanda, S.Si., M.Kom.}%
    };

  \node[anchor=north west, text width=8.95cm, align=left, inner sep=0pt]
    at ($(current page.north east)+(-9.40cm,-0.92cm)$) {%
      {\color{ubAccent}\bfseries\fontsize{8.2}{10.2}\selectfont
       MATERI #1\ \ \textperiodcentered\ \ MINGGU #2}\\[8pt]
      {\color{ubNavy}\bfseries\fontsize{19.4}{23.4}\selectfont
       \hyphenpenalty=10000 \exhyphenpenalty=10000 \raggedright #3}%
    };

  \node[anchor=north west, text width=8.95cm, align=left, inner sep=0pt]
    at ($(current page.north east)+(-9.40cm,-4.30cm)$) {%
      \renewcommand{\arraystretch}{1.30}%
      \begin{tabular}{@{}>{\color{ubMuted}\fontsize{7.3}{9.0}\selectfont}p{1.85cm}
                        >{\color{ubText}\fontsize{8.1}{10.0}\selectfont}p{6.7cm}@{}}
        Sub-CPMK      & #4 \\
        Teknologi     & #5 \\
        Sumber materi & #6 \\
        Blok          & #7 \\
      \end{tabular}%
    };
\end{tikzpicture}
\end{frame}%
}

% --------------------------------------------------------- pembatas bagian
\newcommand{\sectionframe}[3]{%
\begin{frame}[plain]
\begin{tikzpicture}[remember picture, overlay]
  \fill[ubNavy] (current page.south west) rectangle (current page.north east);
  \fill[ubAccent] (current page.north west) rectangle ([xshift=0.10cm]current page.south west);
  \node[anchor=west, inner sep=0pt]
    at ($(current page.west)+(1.00cm,1.45cm)$)
    {\color{ubNavy!52!white}\bfseries\fontsize{58}{58}\selectfont #1};
  \node[anchor=west, inner sep=0pt, text width=11.6cm, align=left]
    at ($(current page.west)+(1.05cm,-0.55cm)$) {%
      {\color{ubAccent}\bfseries\fontsize{8.0}{10.0}\selectfont BAGIAN}\\[3pt]
      {\color{white}\bfseries\fontsize{20.5}{24.5}\selectfont #2}\\[5pt]
      {\color{ubHead}\small #3}%
    };
  \node[anchor=south east, inner sep=0pt]
    at ($(current page.south east)+(-0.60cm,0.32cm)$)
    {\color{ubHead}\fontsize{7.4}{9}\selectfont #1 / 6
     \ \ \textperiodcentered\ \ \insertframenumber};
\end{tikzpicture}
\end{frame}%
}

% --------------------------------------------------------- peta pertemuan
\newcommand{\blokpeta}[6]{%
  \ifnum#2=#3
    \node[rounded corners=3pt, fill=ubNavy, draw=ubNavy, line width=0.5pt,
          text width=3.02cm, minimum height=1.62cm, align=left, inner sep=5pt]
      at (#1,0) {%
        {\color{ubAccent}\bfseries\fontsize{6.8}{8}\selectfont #4}\\[1pt]
        {\color{white}\bfseries\fontsize{7.6}{9}\selectfont #5}\\[2pt]
        {\color{ubHead}\fontsize{6.4}{7.6}\selectfont #6}};
  \else
    \node[rounded corners=3pt, fill=ubSoft, draw=ubHead, line width=0.5pt,
          text width=3.02cm, minimum height=1.62cm, align=left, inner sep=5pt]
      at (#1,0) {%
        {\color{ubAccent}\bfseries\fontsize{6.8}{8}\selectfont #4}\\[1pt]
        {\color{ubNavy}\bfseries\fontsize{7.6}{9}\selectfont #5}\\[2pt]
        {\color{ubMuted}\fontsize{6.4}{7.6}\selectfont #6}};
  \fi
}

% \petapertemuan{blok aktif}{sebelumnya}{sesudahnya}{keluaran}
% \petapertemuan{blok aktif}{sebelumnya}{pertemuan ini}{sesudahnya}{keluaran}
\newcommand{\petapertemuan}[5]{%
\begin{frame}{Peta Pertemuan}
{\scriptsize\color{ubMuted}Kedudukan materi dalam empat blok perkuliahan}\par
\vspace{0.10cm}
\centering
\begin{tikzpicture}
  \blokpeta{-5.58}{1}{#1}{BLOK I}{Materi 1--5}{Game loop, vektor,\\fisika, collision}
  \blokpeta{-1.86}{2}{#1}{BLOK II}{Materi 6--9}{AI NPC: chase, FSM,\\pathfinding A*}
  \blokpeta{1.86}{3}{#1}{BLOK III}{Materi 10--11}{Unity: scene,\\level, physics}
  \blokpeta{5.58}{4}{#1}{BLOK IV}{Materi 12--14}{Aset visual, audio,\\UI, playtest}
\end{tikzpicture}
\par\vspace{0.30cm}
\begin{columns}[T,onlytextwidth]
  \begin{column}{0.315\textwidth}
    \begin{beamercolorbox}[wd=\linewidth,sep=4pt,rounded=true]{block body}
      {\color{ubMuted}\fontsize{6.6}{8}\selectfont SEBELUMNYA}\\[2pt]
      {\footnotesize #2}
    \end{beamercolorbox}
  \end{column}
  \begin{column}{0.315\textwidth}
    \begin{beamercolorbox}[wd=\linewidth,sep=4pt,rounded=true]{block title}
      {\color{ubHead}\fontsize{6.6}{8}\selectfont PERTEMUAN INI}\\[2pt]
      {\footnotesize #3}
    \end{beamercolorbox}
  \end{column}
  \begin{column}{0.315\textwidth}
    \begin{beamercolorbox}[wd=\linewidth,sep=4pt,rounded=true]{block body}
      {\color{ubMuted}\fontsize{6.6}{8}\selectfont SESUDAHNYA}\\[2pt]
      {\footnotesize #4}
    \end{beamercolorbox}
  \end{column}
\end{columns}
\par\vspace{0.26cm}
{\footnotesize\color{ubText}\textbf{Keluaran pertemuan.} #5\par}
\end{frame}%
}

% --------------------------------------------------------- penutup
\newcommand{\penutup}[2]{%
\begin{frame}[plain]
\begin{tikzpicture}[remember picture, overlay]
  \fill[ubNavy] (current page.south west) rectangle (current page.north east);
  \fill[ubAccent] (current page.north west) rectangle ([xshift=0.10cm]current page.south west);

  \node[anchor=north west, text width=13.6cm, align=left, inner sep=0pt]
    at ($(current page.north west)+(1.00cm,-0.95cm)$) {%
      {\color{ubAccent}\bfseries\fontsize{8.2}{10}\selectfont ASESMEN MINGGU INI}\\[6pt]
      {\color{white}\bfseries\fontsize{24}{28}\selectfont #1}%
    };

  \node[anchor=north west, text width=13.6cm, align=left, inner sep=0pt]
    at ($(current page.north west)+(1.00cm,-4.55cm)$) {%
      {\color{ubNavy!0!ubHead}\rule{13.4cm}{0.4pt}}\\[7pt]
      {\color{ubAccent}\bfseries\fontsize{7.4}{9}\selectfont
       PENGAIT KE MATERI BERIKUTNYA}\\[4pt]
      {\color{ubHead}\fontsize{10.6}{13.2}\selectfont #2}%
    };

  \node[anchor=south east, inner sep=0pt]
    at ($(current page.south east)+(-0.60cm,0.32cm)$)
    {\color{ubHead}\fontsize{7.4}{9}\selectfont \insertframenumber};
\end{tikzpicture}
\end{frame}%
}



\renewcommand{\materino}{06}

\begin{document}

\deckcover{06}{6}
  {Pengantar AI NPC dan Perilaku Chase}
  {2.1}
  {Python / Pygame}
  {Bab VIII buku ajar}
  {Blok II \textperiodcentered\ Materi 6--9}

% --------------------------------------------------------- capaian
\begin{frame}{Capaian Pertemuan}
\setbeamerfont{itemize/enumerate body}{size=\scriptsize}
\awalan{Sub-CPMK dan indikator dari RPS, serta tujuan Bab VIII buku ajar tentang
perilaku mengejar.}
\begin{columns}[T,onlytextwidth]
  \begin{column}{0.52\textwidth}
    \begin{beamercolorbox}[wd=\linewidth,sep=5pt,rounded=true]{block title}
      {\fontsize{6.8}{8}\selectfont\color{ubHead}SUB-CPMK 2.1}\\[2pt]
      {\footnotesize Merancang perilaku \emph{NPC} sederhana (misalnya
       \emph{chase}) berdasarkan posisi pemain.}
    \end{beamercolorbox}
    \vspace{0.16cm}

    {\fontsize{7}{8.5}\selectfont\bfseries\color{ubAccent}INDIKATOR}\par
    \vspace{0.06cm}
    \begin{itemize}
      \item Menghitung \emph{vektor arah} NPC menuju pemain dan
            menormalisasinya.
      \item Membatasi besar kecepatan agar gerak NPC konsisten.
      \item Menjelaskan perbedaan \emph{AI reaktif} dan
            \emph{AI berbasis perencanaan}.
    \end{itemize}
  \end{column}
  \begin{column}{0.45\textwidth}
    {\fontsize{7}{8.5}\selectfont\bfseries\color{ubAccent}TUJUAN BAB VIII}\par
    \vspace{0.08cm}
    \begin{itemize}
      \item Menjelaskan susunan keputusan agen: persepsi, keputusan, dan
            tindakan.
      \item Menghitung vektor arah dan jarak dari NPC menuju pemain.
      \item Menormalkan vektor arah lalu mengalikannya dengan kecepatan.
      \item Menjelaskan mengapa vektor jarak langsung sebagai kecepatan itu
            salah.
      \item Menentukan jari-jari persepsi dan jarak henti beserta alasannya.
      \item Membedakan AI reaktif dan AI berbasis perencanaan beserta
            kelemahannya.
    \end{itemize}
  \end{column}
\end{columns}
\end{frame}

% --------------------------------------------------------- peta
\petapertemuan{2}
  {Materi 5 \textperiodcentered\ Sprite, Group, dan Collision Detection}
  {Materi 6 \textperiodcentered\ Pengantar AI NPC dan Perilaku Chase}
  {Materi 7 \textperiodcentered\ Finite State Machine untuk NPC}
  {NPC mengejar pemain di ruang terbuka; garis arah, lingkaran persepsi, dan
   warna NPC memperlihatkan keputusannya; kecepatan terukur tetap pada berbagai
   jarak.}

% ========================================================= bagian 1
\section{Dari yang Digerakkan Menjadi yang Bergerak Sendiri}
\sectionframe{01}{Dari yang Digerakkan Menjadi yang Bergerak Sendiri}
  {Dari benda pasif menjadi agen yang mengamati, memutuskan, dan bertindak}

\begin{frame}{Tiga Asal Kecepatan yang Sudah Dikenal}
\awalan{Sampai bab sebelumnya, kecepatan setiap benda ditentukan dari luar
dirinya. Bab ini mengubah asal kecepatan itu.}
\begin{itemize}
  \item \textbf{Bab V: tombol.} Kecepatan pemain berasal dari tombol yang
        ditekan. Bendanya pasif: ia bergerak hanya karena diperintah.
  \item \textbf{Bab VI: gravitasi.} Kecepatan tegak berubah karena percepatan
        gravitasi. Perubahannya tetap dan sama untuk semua benda; tidak ada
        keputusan di dalamnya.
  \item \textbf{Bab VII: tabrakan.} Kecepatan berubah karena tabrakan.
        Perubahan itu juga bukan keputusan, melainkan akibat dari keadaan.
\end{itemize}
\vspace{0.16cm}
\begin{kotak}
Pada bab ini kecepatan berasal dari \textbf{keputusan benda itu sendiri}
berdasarkan keadaan di sekitarnya. Bendanya menjadi agen, yaitu sesuatu yang
mengamati, memutuskan, dan bertindak. Pandangan permainan pun berganti menjadi
pandangan atas, sehingga gravitasi tidak diperlukan.
\end{kotak}
\end{frame}

\begin{frame}{Susunan Keputusan Sebuah Agen}
\awalan{Tiga langkah berikut selalu ada pada agen mana pun, dari yang paling
sederhana sampai yang paling rumit.}
\begin{enumerate}
  \item \textbf{Persepsi.} Mengumpulkan keterangan tentang keadaan di
        sekitarnya. Pada bab ini keterangan itu hanya satu: di mana pemain
        berada sekarang.
  \item \textbf{Keputusan.} Menentukan tindakan yang dipilih. Pada bab ini
        hanya ada dua pilihan, yaitu mengejar atau berdiam.
  \item \textbf{Tindakan.} Melaksanakan keputusan itu dengan mengubah
        kecepatan, lalu posisi.
\end{enumerate}
\vspace{0.16cm}
\begin{kotak}
Keluwesan susunan itu yang menarik: menambahkan perilaku baru berarti
menambahkan pilihan pada langkah kedua, sedangkan langkah pertama dan ketiga
tidak perlu berubah. Sifat itulah yang membuat susunan ini dapat dikembangkan
menjadi mesin keadaan pada Materi 7 dan dilengkapi pencarian jalur pada
Materi 8.
\end{kotak}
\end{frame}

\begin{frame}{Perilaku yang Dibangun dan yang Sengaja Belum Ada}
\awalan{Perilaku yang dibangun sederhana dan sengaja demikian: NPC mengejar
apabila pemain di dalam jangkauannya, dan berdiam apabila pemain di luar
jangkauan.}
\begin{itemize}
  \item \textbf{Keadaan lebih dari dua}, misalnya berpatroli dan menyerang.
        Dibahas pada Materi 7 melalui mesin keadaan.
  \item \textbf{Pengetahuan tentang rintangan}, sehingga NPC dapat mencari
        jalan memutar. Dibahas pada Materi 8 melalui pencarian jalur A*.
  \item \textbf{Pertimbangan biaya perhitungan} ketika jumlah NPC bertambah.
        Dibahas pada Materi 9.
\end{itemize}
\vspace{0.16cm}
\begin{kotak}
Walaupun sederhana, perilaku ini sudah memuat seluruh persoalan penting:
bagaimana jarak dihitung, bagaimana arah ditentukan, bagaimana kecepatan
dibatasi, dan kapan berhenti.
\end{kotak}
\end{frame}

% ========================================================= bagian 2
\section{Persepsi Jarak dan Arah}
\sectionframe{02}{Persepsi Jarak dan Arah}
  {Satu pengurangan menghasilkan dua keterangan yang berbeda perannya}

\begin{frame}{Jarak untuk Memutuskan, Arah untuk Bergerak}
\awalan{Persepsi adalah langkah pertama. Yang diamati hanya posisi pemain, dan
yang dihitung dari keterangan itu ada dua: jarak dan arah.}
\begin{center}
\adjustbox{max width=0.92\linewidth,max height=0.53\textheight}{%
  % chapters/08/figures/persepsi-npc.tikz
% Persepsi NPC: jarak ke pemain menentukan keputusan, arah menuju pemain
% menentukan kecepatan, dan besarnya kecepatan tetap berapa pun jaraknya.
\begin{tikzpicture}[
  npc/.style={rectangle, draw=ubRed, thick, fill=ubRed!25, inner sep=4pt},
  pemain/.style={rectangle, draw=ubPrimary, thick, fill=ubPrimary!25, inner sep=4pt},
  lingkar/.style={draw=ubLine, thick, dashed},
  lingkarkecil/.style={draw=ubGreen, thick},
  garisjarak/.style={draw=ubWarn, thick, dashed},
  panah/.style={-{Stealth[length=2.4mm]}, very thick, draw=ubAccent},
  label/.style={font=\scriptsize, align=left}
]

% NPC di titik asal, dengan lingkaran persepsi dan lingkaran jarak henti.
\draw[lingkar] (0,0) circle (3.0);
\draw[lingkarkecil] (0,0) circle (0.55);
\node[npc] (npc) at (0,0) {NPC};

% Pemain berada di dalam jari-jari persepsi.
\node[pemain] (player) at (2.3,1.8) {P};

% Garis jarak dan garis arah.
\draw[garisjarak] (npc.center) -- (player.center);
\node[label, ubWarn, anchor=south east] at (1.35,1.15) {$d$};

\draw[panah] (0.3,0.24) -- (1.15,0.9);
\node[label, ubAccent, anchor=west] at (1.05,0.72)
  {kecepatan $= \dfrac{\text{arah}}{|d|} \times v$};

% Keterangan lingkaran.
\node[label, ubLine, anchor=west] at (2.15,2.15) {jari-jari persepsi};
\draw[ubLine, thin] (2.1,2.1) -- (2.0,1.35);

\node[label, ubGreen, anchor=north west] at (0.5,-0.75) {jarak henti};
\draw[ubGreen, thin] (0.55,-0.7) -- (0.38,-0.42);

% Keterangan perilaku.
\node[label, anchor=north, text width=104mm] at (0,-3.5)
  {Apabila pemain berada di dalam jari-jari persepsi dan di luar jarak henti,
   NPC bergerak menuju pemain dengan kecepatan tetap. Perhatikan bahwa panjang
   panah kecepatan tidak bergantung pada jarak $d$.};
\end{tikzpicture}
}
\end{center}
\vspace{-0.02cm}
\begin{itemize}
  \item Jarak dan arah diperoleh dari satu perhitungan yang sama, yaitu selisih
        kedua posisi: $\vec{u} = \vec{p} - \vec{n}$.
  \item Jarak dipakai untuk memutuskan, sedangkan vektor itu dipakai untuk
        bergerak.
\end{itemize}
\end{frame}

\begin{frame}[fragile]{Satu Pengurangan, Dua Hasil}
\awalan{Dua baris pertama \texttt{Npc.update} menghasilkan jarak dan arah
sekaligus.}
\lstinputlisting[language=Python,firstline=124,lastline=125]
  {../buku_ajar/code/python/ch08_chase.py}
\kodejelas{Baris pertama menghasilkan vektor yang menunjuk dari NPC ke pemain.
Perhatikan urutan pengurangannya: yang dikurangi adalah posisi NPC, sehingga
vektornya menunjuk \textbf{dari NPC menuju pemain}. Baris kedua menghitung
panjang vektor itu, yaitu jaraknya dalam piksel.}
\vspace{0.15cm}
\begin{peringatan}
Bila urutannya dibalik, vektornya menunjuk arah sebaliknya dan NPC justru
menjauh. Kekeliruan itu sering terjadi karena kedua urutan sama-sama masuk
akal ketika ditulis.
\end{peringatan}
\end{frame}

\begin{frame}{Jari-jari Persepsi sebagai Batas Keputusan}
\awalan{NPC tidak seharusnya mengetahui keberadaan pemain di mana pun.
Jangkauan penglihatannya dibatasi sebuah nilai yang disebut jari-jari
persepsi.}
\begin{itemize}
  \item Nilai itu dibandingkan dengan jarak: bila jarak tidak lebih besar
        daripada jari-jari, pemain terlihat dan keputusan berikutnya adalah
        mengejar.
  \item Nilainya urusan perancangan, bukan urusan teknis. Jari-jari terlalu
        besar membuat NPC terasa mengetahui segala hal; jari-jari terlalu kecil
        membuat NPC mudah dihindari sehingga tantangan hilang.
  \item Pada program bab ini nilainya dapat diubah melalui baris perintah agar
        pengaruhnya dapat dirasakan langsung.
  \item Program menggambar empat hal agar keputusan terlihat: lingkaran
        jari-jari persepsi, lingkaran jarak henti, garis arah menuju pemain,
        dan warna NPC yang berubah menurut keputusannya.
\end{itemize}
\vspace{0.12cm}
\begin{kotak}
Tanpa gambar itu, Anda hanya melihat NPC yang bergerak atau tidak bergerak,
dan tidak dapat membedakan ``berdiam karena memutuskan berdiam'' dari
``berdiam karena kesalahan perhitungan''.
\end{kotak}
\end{frame}

% ========================================================= bagian 3
\section{Menormalkan Arah dan Membatasi Kecepatan}
\sectionframe{03}{Menormalkan Arah dan Membatasi Kecepatan}
  {Memisahkan arah dari besar kecepatan}

\begin{frame}{Dua Cara Mengubah Arah Menjadi Kecepatan}
\awalan{Dengan $\vec{u} = \vec{p} - \vec{n}$ sebagai selisih posisi pemain dan
NPC, serta $s$ sebagai kecepatan yang Anda tetapkan, ada dua cara yang perlu
dibandingkan.}
\begin{columns}[T,onlytextwidth]
  \begin{column}{0.47\textwidth}
    \begin{block}{Cara yang salah}
      \[
        \vec{v} = \vec{u},
        \qquad \lVert \vec{v} \rVert = d
      \]
      {\scriptsize Arahnya benar, tetapi besarnya ikut berubah menurut jarak:
       pada 400 piksel kecepatannya 400 piksel per detik; pada 75 piksel,
       75 piksel per detik.}
    \end{block}
  \end{column}
  \begin{column}{0.49\textwidth}
    \begin{exampleblock}{Cara yang benar}
      \[
        \hat{u} = \frac{\vec{u}}{\lVert \vec{u} \rVert},
        \qquad \vec{v} = \hat{u}\, s
      \]
      {\scriptsize Tidak ada lagi nilai jarak pada perhitungan itu. Yang
       tersisa hanyalah arah dan kecepatan yang Anda tetapkan.}
    \end{exampleblock}
  \end{column}
\end{columns}
\vspace{0.10cm}
\begin{kotak}
\emph{Normalisasi} memisahkan arah dari besar kecepatan. Arah berasal dari
keadaan permainan, sedangkan besarnya berasal dari rancangan Anda.
\end{kotak}
\end{frame}

\begin{frame}{Mengapa Kecepatan yang Tetap Itu Penting}
\awalan{Kecepatan yang bergantung pada jarak bukan sekadar kurang rapi; ia
merusak permainan dengan tiga cara.}
\begin{enumerate}
  \item \textbf{Perilaku NPC tidak dapat diramalkan pemain.} Pemain tidak dapat
        belajar kapan ia aman dan kapan tidak, sebab kecepatan pengejar berubah
        menurut jaraknya. Pada permainan yang adil, pemain harus dapat
        memperkirakan akibat dari tindakannya.
  \item \textbf{Tingkat kesulitan berubah tanpa disengaja.} Ruang yang luas
        membuat pemain lebih mudah lolos, sedangkan ruang yang sempit membuatnya
        lebih sulit. Tingkat kesulitan menjadi bergantung pada ukuran ruang,
        bukan pada rancangan Anda.
  \item \textbf{Penghitungan waktu menjadi keliru.} Perkiraan bahwa NPC akan
        mencapai suatu tempat dalam dua detik hanya benar pada satu jarak
        tertentu.
\end{enumerate}
\vspace{0.12cm}
\begin{kotak}
Sebagai pedoman yang berlaku umum: \textbf{arah dan besar kecepatan sebaiknya
ditentukan secara terpisah}. Pedoman itu berlaku juga untuk berpatroli,
menghindar, dan seluruh perilaku lain pada bab-bab berikutnya.
\end{kotak}
\end{frame}

\begin{frame}{Perbandingan yang Dapat Diukur}
\awalan{Perbedaan kedua cara dibuktikan dengan mengukur kecepatan NPC pada
berbagai jarak, dengan kecepatan ditetapkan 150 piksel per detik.}
{\scriptsize
\renewcommand{\arraystretch}{1.16}
\begin{tabularx}{\textwidth}{@{}>{\raggedright\arraybackslash}p{3.7cm}
    >{\raggedright\arraybackslash}p{2.5cm}
    >{\raggedright\arraybackslash}p{3.5cm}
    >{\raggedright\arraybackslash}X@{}}
\toprule
\textbf{Cara} & \textbf{Jarak saat diukur} & \textbf{Kecepatan terukur} &
\textbf{Keterangan} \\
\midrule
Arah dinormalkan & 59--167 piksel & 150,0 piksel/detik &
  Tetap, tidak bergantung jarak. \\
\addlinespace
Vektor jarak dipakai langsung & 75--184 piksel & 75,2--184,0 piksel/detik &
  Kecepatan sama persis dengan jaraknya. \\
\bottomrule
\end{tabularx}}
\vspace{0.18cm}
\begin{kotak}
Perhatikan baris kedua: kecepatannya sama persis dengan jaraknya. Itulah tanda
khas cara yang salah. Pada baris pertama, kecepatannya tetap walaupun jaraknya
berubah hampir tiga kali lipat.
\end{kotak}
\end{frame}

% ========================================================= bagian 4
\section{Kapan Mengejar dan Kapan Berhenti}
\sectionframe{04}{Kapan Mengejar dan Kapan Berhenti}
  {Dua batas, satu getaran yang harus dicegah}

\begin{frame}{Getaran di Atas Pemain}
\awalan{Perilaku mengejar memerlukan dua batas, bukan satu. Batas pertama
adalah jari-jari persepsi; batas kedua adalah jarak henti.}
\begin{itemize}
  \item Bila NPC terus mengejar sampai jaraknya nol, ia tidak dapat berhenti
        tepat pada posisi pemain: ia melampauinya, berbalik, lalu melampauinya
        lagi.
  \item Pada kecepatan 150 piksel per detik dan laju 60 bingkai per detik,
        perpindahan setiap bingkainya 2,5 piksel. NPC bergetar di sekitar
        pemain, tampak menempel padanya, dan mengganggu pemeriksaan tabrakan.
\end{itemize}
\[
  \Delta = \frac{v}{\text{laju bingkai}} = \frac{150}{60} = 2{,}5\
  \text{piksel},
  \qquad d_{\text{henti}} > \Delta
\]
\begin{kiat}
Agar getaran tidak terjadi, jarak henti harus lebih besar daripada perpindahan
satu bingkai. Bila kecepatan dinaikkan menjadi 3.000 piksel per detik,
perpindahannya menjadi 50 piksel dan melampaui jarak henti, sehingga NPC
melewati daerah henti tanpa pernah berhenti di dalamnya.
\end{kiat}
\end{frame}

\begin{frame}[fragile]{Urutan Pemeriksaan Jarak}
\awalan{Jarak henti diperiksa lebih dahulu, baru jari-jari persepsi, sebab jarak
henti selalu lebih kecil daripada jari-jari.}
\lstinputlisting[language=Python,firstline=127,lastline=130]
  {../buku_ajar/code/python/ch08_chase.py}
\kodejelas{Bila jarak sudah tidak lebih besar daripada jarak henti, kecepatan
dinolkan dan penanda \texttt{chasing} dimatikan. Cabang berikutnya baru
menangani pengejaran. Membalik urutan ini membuat cabang jari-jari menelan
daerah henti, sehingga NPC tidak pernah berhenti.}
\vspace{0.15cm}
\begin{kotak}
Pemeriksaan dilakukan berurutan, bukan sekaligus: cabang pertama yang cocok
yang menentukan keputusan pada bingkai itu.
\end{kotak}
\end{frame}

\begin{frame}{Hysteresis pada Batas Persepsi}
\awalan{Persoalan kecil muncul tepat pada batas jari-jari persepsi. Bila pemain
berjalan mengelilingi NPC pada jarak itu, keputusan NPC berganti-ganti pada
setiap bingkai karena jaraknya naik turun di sekitar nilai batas.}
\begin{itemize}
  \item Gejalanya terlihat sebagai NPC yang bergetar dan berkedip warnanya.
  \item Penanganannya: pakai \textbf{dua nilai batas}, bukan satu. NPC mulai
        mengejar ketika jaraknya kurang dari jari-jari persepsi, tetapi tidak
        berhenti mengejar sebelum jaraknya melampaui nilai yang sedikit lebih
        besar.
  \item Cara itu disebut \emph{hysteresis}, dan dipakai pada banyak alat
        pengatur untuk mencegah pergantian keadaan yang berlebihan.
  \item Variabel \texttt{self.chasing} mengingat keputusan pada bingkai
        sebelumnya, sehingga keputusan kini bergantung juga pada keadaan
        sebelumnya. Itulah keadaan pertama yang dimiliki sebuah NPC.
\end{itemize}
\vspace{0.12cm}
\begin{kotak}
Menambah keadaan lain, misalnya berpatroli dan menyerang, beserta aturan
perpindahan di antaranya, akan menghasilkan mesin keadaan yang dibahas pada
Materi 7.
\end{kotak}
\end{frame}

% ========================================================= bagian 5
\section{Program Perilaku Mengejar}
\sectionframe{05}{Program Perilaku Mengejar}
  {Tiga langkah keputusan di dalam satu fungsi}

\begin{frame}{Seluruh Keputusan pada Satu Fungsi}
\awalan{Seluruh keputusan NPC terletak pada satu fungsi, yaitu
\texttt{Npc.update}. Susunannya mengikuti tiga langkah pada bagian pertama.}
\begin{itemize}
  \item \textbf{Persepsi.} Dua baris pertama menghitung vektor menuju pemain
        dan jaraknya. Keduanya dihitung dari satu pengurangan yang sama,
        sehingga tidak ada perhitungan yang berulang.
  \item \textbf{Keputusan.} Rangkaian pemeriksaan jarak menentukan kecepatan.
        Urutannya penting: jarak henti diperiksa lebih dahulu, baru jari-jari
        persepsi.
  \item \textbf{Tindakan.} Posisi diperbarui dengan kecepatan tersebut,
        dibatasi agar tidak keluar jendela, lalu \texttt{rect} diselaraskan
        sebagaimana biasa.
\end{itemize}
\vspace{0.16cm}
\begin{kotak}
Kelas \texttt{Npc} tidak menerima arah dari luar, berbeda dari kelas
\texttt{Player} pada bab ini dan pada bab-bab sebelumnya. Itulah perbedaan yang
paling mendasar antara benda yang diperintah dan benda yang memutuskan.
\end{kotak}
\end{frame}

\begin{frame}[fragile]{Arah Dinormalkan, Besar Kecepatan Tetap}
\awalan{Cabang pengejaran memuat kedua cara perhitungan; pilihan
\texttt{--naive} menentukan cabang mana yang dipakai.}
\lstinputlisting[language=Python,firstline=133,lastline=141]
  {../buku_ajar/code/python/ch08_chase.py}
\kodejelas{Cabang \texttt{naive} memakai vektor jarak langsung, sehingga
kecepatan sama dengan jarak. Cabang yang benar menormalkan vektor lebih dahulu,
lalu mengalikannya dengan \texttt{self.speed}, sehingga besarnya tetap berapa
pun jaraknya.}
\vspace{0.12cm}
\begin{kotak}
Jalankan program dengan dan tanpa \texttt{--naive}, lalu bandingkan kecepatan
terukur pada terminal. Perbedaan itulah bukti indikator (b) pada bab ini.
\end{kotak}
\end{frame}

\begin{frame}{Percobaan yang Disarankan}
\awalan{Setiap nilai yang Anda tetapkan dapat diubah dari baris perintah, dan
pengaruhnya dapat langsung diamati.}
{\scriptsize
\renewcommand{\arraystretch}{1.10}
\begin{tabularx}{\textwidth}{@{}>{\raggedright\arraybackslash}p{4.6cm}
    >{\raggedright\arraybackslash}X@{}}
\toprule
\textbf{Yang diubah dan caranya} & \textbf{Yang sebaiknya Anda amati} \\
\midrule
Jari-jari persepsi, melalui \texttt{--radius}. &
  Pada jari-jari kecil, pemain dapat berlalu-lalang dekat NPC tanpa dikejar.
  Pada jari-jari besar, NPC mulai bergerak bahkan ketika pemain jauh. \\
\addlinespace
Kecepatan NPC, melalui \texttt{--speed}. &
  Bila jauh lebih rendah daripada pemain, NPC tidak pernah menangkap. Bila
  lebih tinggi, pemain selalu tertangkap dan NPC berhenti pada jarak
  hentinya. \\
\addlinespace
Jarak henti, dengan mengubah \texttt{DEFAULT\_STOP} menjadi 2. &
  NPC bergetar di sekitar pemain, karena perpindahan setiap bingkainya lebih
  besar daripada jarak henti. \\
\addlinespace
Cara menghitung kecepatan, melalui \texttt{--naive}. &
  Kecepatan berubah menurut jarak: lambat ketika jauh, cepat ketika dekat. \\
\addlinespace
Banyaknya NPC, melalui \texttt{--npc}. &
  Gerombolan mengelilingi pemain; perhatikan pula beban perhitungan ketika
  jumlahnya bertambah besar. \\
\bottomrule
\end{tabularx}}
\vspace{0.12cm}
\begin{kotak}
Pilihan \texttt{--naive} sengaja disediakan untuk memperlihatkan cara yang
salah. Perbedaan kecepatan terukur itulah bukti indikator (b).
\end{kotak}
\end{frame}

% ========================================================= bagian 6
\section{AI Reaktif dan AI Berbasis Perencanaan}
\sectionframe{06}{AI Reaktif dan AI Berbasis Perencanaan}
  {Kecukupan, bukan mana yang lebih baik}

\begin{frame}{Ketika Ada Rintangan di Antara Mereka}
\awalan{NPC pada bab ini disebut reaktif karena seluruh keputusannya diambil
dari keadaan yang teramati saat itu.}
\begin{center}
\adjustbox{max width=0.86\linewidth,max height=0.43\textheight}{%
  % chapters/08/figures/reaktif-vs-perencanaan.tikz
% Dua cara menyusun keputusan NPC: reaktif yang hanya melihat arah pemain,
% dan berbasis perencanaan yang memperhitungkan rintangan.
\begin{tikzpicture}[
  npc/.style={rectangle, draw=ubRed, thick, fill=ubRed!25, inner sep=4pt},
  pemain/.style={rectangle, draw=ubPrimary, thick, fill=ubPrimary!25, inner sep=4pt},
  dinding/.style={rectangle, draw=ubLine, thick, fill=ubGray},
  panah/.style={-{Stealth[length=2.2mm]}, very thick, draw=ubRed},
  panahjalan/.style={-{Stealth[length=2.2mm]}, very thick, draw=ubGreen},
  judul/.style={font=\footnotesize\bfseries, ubPrimary, align=center},
  label/.style={font=\scriptsize, align=center}
]

% ---------- Panel A: reaktif ----------
\begin{scope}
  \node[judul, anchor=south] at (2.4,2.35) {AI reaktif};
  \draw[dinding] (2.3,0.5) rectangle (2.65,2.3);

  \node[npc] (n1) at (0.8,1.4) {NPC};
  \node[pemain] (p1) at (4.0,1.4) {P};

  \draw[panah] (1.15,1.4) -- (2.2,1.4);
  \node[label, ubRed, anchor=north, text width=44mm] at (2.4,0.2)
    {bergerak lurus ke arah pemain, sehingga terus menabrak dinding dan tidak
     pernah sampai};
\end{scope}

% ---------- Panel B: berbasis perencanaan ----------
\begin{scope}[xshift=6.4cm]
  \node[judul, anchor=south] at (2.4,2.35) {AI berbasis perencanaan};
  \draw[dinding] (2.3,0.5) rectangle (2.65,2.3);

  \node[npc] (n2) at (0.8,1.4) {NPC};
  \node[pemain] (p2) at (4.0,1.4) {P};

  \draw[panahjalan] (1.15,1.4) -- (1.5,1.4);
  \draw[panahjalan] (1.5,1.4) -- (1.5,0.15) -- (3.5,0.15);
  \draw[panahjalan] (3.5,0.15) -- (3.5,1.4) -- (3.68,1.4);

  \node[label, ubGreen, anchor=north, text width=46mm] at (2.4,0.2)
    {memperhitungkan rintangan lebih dahulu, sehingga menemukan jalan
     memutar};
\end{scope}
\end{tikzpicture}
}
\end{center}
\vspace{-0.02cm}
\begin{itemize}
  \item AI reaktif hanya mengetahui arah menuju pemain. Dinding tidak termasuk
        dalam persepsinya, sehingga NPC bergerak lurus dan menekan dinding
        selama pemain berada di baliknya.
  \item AI berbasis perencanaan memeriksa apakah jalan lurusnya terhalang,
        lalu mencari jalan memutar.
\end{itemize}
\end{frame}

\begin{frame}{Kelebihan dan Kelemahan AI Reaktif}
\awalan{Kesederhanaan AI reaktif membawa tiga keuntungan, tetapi juga satu
kelemahan yang langsung terasa.}
\begin{columns}[T,onlytextwidth]
  \begin{column}{0.47\textwidth}
    \begin{exampleblock}{Keuntungan}
      \begin{itemize}
        \item Kodenya pendek dan mudah diperiksa.
        \item Perhitungannya murah, sehingga jumlah NPC dapat banyak.
        \item Perilakunya mudah diramalkan oleh perancang, sehingga mudah
              diatur.
      \end{itemize}
    \end{exampleblock}
  \end{column}
  \begin{column}{0.49\textwidth}
    \begin{block}{Kelemahan}
      \begin{itemize}
        \item Begitu ada rintangan di antara NPC dan pemain, NPC tidak
              mengetahuinya.
        \item Ia bergerak lurus ke arah pemain, menabrak dinding, dan tetap
              menekan dinding itu.
        \item Dari sisi pemain, NPC tampak bodoh: ia berhenti di depan dinding
              tanpa berusaha mencari jalan lain.
      \end{itemize}
    \end{block}
  \end{column}
\end{columns}
\vspace{0.12cm}
\begin{kotak}
NPC itu tidak melakukan kesalahan. Ia melaksanakan keputusannya dengan benar.
Yang kurang adalah \textbf{\mbox{keterangan}}---data yang tersedia baginya
untuk mengambil keputusan.
\end{kotak}
\end{frame}

\begin{frame}{Kapan Masing-Masing Dipakai}
\awalan{Pilihan di antara keduanya bukan soal mana yang lebih baik, melainkan
soal kecukupan.}
{\scriptsize
\renewcommand{\arraystretch}{1.10}
\begin{tabularx}{\textwidth}{@{}>{\raggedright\arraybackslash}p{4.6cm}
    >{\raggedright\arraybackslash}p{3.4cm}
    >{\raggedright\arraybackslash}X@{}}
\toprule
\textbf{Keadaan permainan} & \textbf{Cara yang memadai} & \textbf{Alasan} \\
\midrule
Ruang terbuka tanpa rintangan & Reaktif &
  Tidak ada yang menghalangi, sehingga arah menuju pemain selalu dapat
  ditempuh. Inilah keadaan pada bab ini. \\
\addlinespace
Rintangan kecil yang mudah dihindari & Reaktif, ditambah penghindaran
  sederhana &
  Perencanaan penuh terlalu mahal untuk rintangan yang jarang. \\
\addlinespace
Labirin atau banyak rintangan & Berbasis perencanaan &
  Jalan lurus hampir selalu terhalang, sehingga diperlukan pencarian rute.
  Dibahas pada Materi 8. \\
\addlinespace
Banyak NPC yang berpindah secara terkoordinasi & Berbasis perencanaan dengan
  penghematan &
  Biaya pencarian jalur menjadi pertimbangan utama. Dibahas pada Materi 9. \\
\bottomrule
\end{tabularx}}
\vspace{0.14cm}
\begin{kotak}
Kedua kekurangan NPC pada bab ini berbeda jenisnya: yang pertama berkaitan
dengan \textbf{kualitas keputusan}, sedangkan yang kedua berkaitan dengan
\textbf{biaya keputusan}.
\end{kotak}
\end{frame}

% --------------------------------------------------------- aktivitas
\begin{frame}{Aktivitas di Kelas \textperiodcentered\ Praktikum 06}
\awalan{Pengalaman belajar minggu 6 pada RPS: NPC mengejar pemain di ruang
terbuka, dengan garis arah yang divisualisasikan agar keputusannya dapat
diamati.}
\begin{enumerate}
  \item Ketik ulang \texttt{ch08\_chase.py} ke folder \texttt{code}, jalankan,
        lalu gerakkan pemain dengan WASD. Amati warna NPC dan garis arah yang
        menunjuk ke pemain.
  \item Amati batas persepsi: berjalanlah menjauh sampai seluruh NPC berhenti,
        lalu dekati kembali. Catat jarak ketika NPC mulai bergerak.
  \item Ukur kecepatan pada berbagai jarak dengan
        \texttt{python ch08\_chase.py --frames 400 --auto-run}, lalu catat
        kecepatan terukur pada terminal.
  \item Jalankan kembali dengan \texttt{--naive}, catat kecepatan terukurnya,
        dan jelaskan perbedaan kedua hasil.
  \item Jalankan program dengan jari-jari 120, 240, dan 400 piksel; catat kesan
        tingkat kesulitan pada masing-masing nilai.
  \item Tambahkan bongkahan tegak di tengah ruang agar pemain dapat
        bersembunyi, catat perilaku NPC, lalu simpan pekerjaan ke repositori
        Git.
\end{enumerate}
\vspace{0.10cm}
{\footnotesize\color{ubAccent}\itshape Rincian lengkap: Praktikum 06 pada
Bab VIII buku ajar. Kumpulkan melalui repositori Git pribadi.}
\end{frame}

% --------------------------------------------------------- rangkuman
\begin{frame}{Rangkuman}
\begin{itemize}
  \item Benda menjadi agen: kecepatannya tidak lagi berasal dari tombol atau
        gravitasi, melainkan dari keputusannya sendiri.
  \item Susunan keputusan agen terdiri atas persepsi, keputusan, dan tindakan;
        susunan itu tetap sama dari agen paling sederhana sampai paling rumit.
  \item Persepsi menghasilkan dua hal dari satu perhitungan: vektor menuju
        pemain dan jaraknya. Jarak dipakai untuk memutuskan, arah untuk
        bergerak.
  \item Kecepatan dihitung dengan menormalkan vektor arah lalu mengalikannya
        dengan kecepatan yang ditetapkan; vektor jarak langsung sebagai
        kecepatan adalah cara yang salah.
  \item Jarak henti mencegah getaran di atas pemain dan harus lebih besar
        daripada perpindahan satu bingkai; \emph{hysteresis} mencegah
        pergantian keputusan pada batas persepsi.
  \item AI reaktif murah dan mudah diperiksa tetapi tidak mengetahui rintangan;
        AI berbasis perencanaan menuntut perhitungan yang lebih besar.
\end{itemize}
\end{frame}

% --------------------------------------------------------- cek
\begin{frame}{Cek Pemahaman}
\begin{enumerate}
  \item Mengapa urutan pengurangan pada perhitungan vektor tidak boleh dibalik,
        dan apa akibatnya bila dibalik?
  \item Mengapa vektor jarak tidak boleh dipakai langsung sebagai kecepatan?
        Sertakan satu contoh angkanya.
  \item Sebuah NPC reaktif menekan dinding selama pemain berada di baliknya.
        Apakah itu kesalahan NPC? Jelaskan.
\end{enumerate}
\vspace{0.35cm}
{\footnotesize\color{ubMuted}\itshape Jawab tanpa membuka catatan, lalu
bandingkan dengan pasangan di sebelah Anda.}
\end{frame}

% --------------------------------------------------------- penutup
\penutup{Tugas Praktikum 06}
  {Materi 7 menyusun \emph{Finite State Machine} untuk NPC: keadaan
   \emph{idle}, \emph{patrol}, dan \emph{chase} beserta aturan perpindahannya,
   sehingga NPC dapat berpatroli dan mengelak, bukan hanya mengejar atau
   berdiam.}

\end{document}
